\chapter{مجاميع المتغيرات العشوائية وقانون الأعداد الكبيرة}\label{ch:g12:sums}

لماذا تربح دور القمار دائمًا في النهاية، ولماذا تنجح استطلاعات الرأي؟ لأن
متوسطات مقادير عشوائية مستقلة كثيرة تتذبذب أقل فأقل. ويبرهن
هذا الفصل على ذلك: خطية \omterm{def:g12:randvar:exp}{الأمل الرياضي}، وجمعية \omterm{def:g12:randvar:exp}{التباين} من أجل
المتغيرات المستقلة، ومتراجحة بيينيمي--تشيبيشيف، وقانون الأعداد
الكبيرة.

\section{مجاميع المتغيرات العشوائية}

إذا أُعطي متغيران عشوائيان $X, Y$ على \omterm{def:g11:prob:model}{مجموعة الإمكانيات} المنتهية نفسها
$\Omega$، فإن المجموع $X + Y$ هو \omterm{def:g12:randvar:rv}{المتغير العشوائي}
$\omega \mapsto X(\omega) + Y(\omega)$.

\begin{theorem}[خطية الأمل الرياضي]\label{thm:g12:sums:linearity}
\index{أمل رياضي!خطية}
من أجل كل متغيرين عشوائيين $X, Y$ على $\Omega$ وكل $a, b \in \R$:
\[
\E(X + Y) = \E(X) + \E(Y), \qquad \E(aX + b) = a\E(X) + b .
\]
\emph{ولا يلزم أي فرض \omterm{def:g12:condprob:indep}{استقلال}.}
\end{theorem}

\begin{proof}
اكتب \omterm{def:g12:randvar:exp}{الأمل الرياضي} مجموعًا على الإمكانيات: فبما أن
$\P(X = x) = \sum_{\omega : X(\omega) = x} \P(\{\omega\})$، يعطي تجميع الحدود
أن $\E(X) = \sum_{\omega \in \Omega} \P(\{\omega\})\,X(\omega)$. ثم
\[
\E(X + Y) = \sum_{\omega} \P(\{\omega\})\bigl(X(\omega) + Y(\omega)\bigr)
= \sum_{\omega} \P(\{\omega\})X(\omega)
+ \sum_{\omega} \P(\{\omega\})Y(\omega) = \E(X) + \E(Y). \qedhere
\]
\end{proof}

\begin{definition}[المتغيران العشوائيان المستقلان]\label{def:g12:sums:indep}
يكون $X$ و $Y$ \emph{مستقلين}\index{متغير عشوائي!مستقل} إذا كان
من أجل كل قيمتين $x, y$:
\[
\P(X = x \text{ و } Y = y) = \P(X = x)\,\P(Y = y).
\]
وتكون عدة متغيرات $X_1, \dots, X_n$ مستقلة إذا صحت قاعدة الجداء
هذه من أجل كل اختيار لقيم كل عائلة جزئية.
\end{definition}

\begin{proposition}\label{prop:g12:sums:prodexp}
إذا كان $X$ و $Y$ \omterm{def:g12:sums:indep}{مستقلين}، فإن $\E(XY) = \E(X)\,\E(Y)$.
\end{proposition}

\begin{proof}
\[
\begin{aligned}
\E(XY) &= \sum_{x, y} xy\;\P(X = x \text{ و } Y = y)
= \sum_{x, y} xy\,\P(X=x)\P(Y=y) \\
&= \Bigl(\sum_x x\P(X=x)\Bigr)\Bigl(\sum_y y\P(Y=y)\Bigr). \qedhere
\end{aligned}
\]
\end{proof}

\begin{theorem}[تباين مجموع]\label{thm:g12:sums:variance}
\index{تباين!مجموع}
إذا كان $X$ و $Y$ \emph{\omterm{def:g12:sums:indep}{مستقلين}}، فإن
\[
\V(X + Y) = \V(X) + \V(Y).
\]
وبوجه أعم، من أجل $X_1, \dots, X_n$ مستقلة:
$\V(X_1 + \dots + X_n) = \V(X_1) + \dots + \V(X_n)$.
\end{theorem}

\begin{proof}
باستعمال صيغة كونيغ--هويغنز (\cref{prop:g12:randvar:konig}) والخطية:
\begin{align*}
\V(X + Y) &= \E\bigl((X+Y)^2\bigr) - \bigl(\E X + \E Y\bigr)^2\\
&= \E(X^2) + 2\E(XY) + \E(Y^2)
- \E(X)^2 - 2\E(X)\E(Y) - \E(Y)^2\\
&= \V(X) + \V(Y) + 2\bigl(\E(XY) - \E(X)\E(Y)\bigr),
\end{align*}
وينعدم القوس الأخير من أجل المتغيرات المستقلة
(\cref{prop:g12:sums:prodexp}). وتنتج الحالة العامة بالتراجع.
\end{proof}

\begin{example}[التوزيع الثنائي من جديد]\label{ex:g12:sums:binomial}
المتغير الثنائي $X \sim \mathcal B(n, p)$ هو مجموع
$X = X_1 + \dots + X_n$ لعدد $n$ من متغيرات برنولي المستقلة. ومنه، بنيويًا:
\[
\E(X) = np, \qquad \V(X) = np(1-p),
\]
فنستعيد \cref{thm:g12:randvar:binomial} بلا أي حساب.
\end{example}

\begin{proposition}[متوسط العينة]\label{prop:g12:sums:mean}
لتكن $X_1, \dots, X_n$ مستقلة ولها \omterm{def:g12:randvar:rv}{التوزيع} نفسه الذي للمتغير $X$
(\omterm{def:g12:randvar:exp}{الأمل الرياضي} $\mu$ \omterm{def:g12:randvar:exp}{والتباين} $\sigma^2$)، وليكن
$M_n = \frac{X_1 + \dots + X_n}{n}$ \emph{متوسط
عينتها}\index{متوسط العينة}. عندئذٍ
\[
\E(M_n) = \mu, \qquad \V(M_n) = \frac{\sigma^2}{n}, \qquad
\sigma(M_n) = \frac{\sigma}{\sqrt n}.
\]
\end{proposition}

\begin{proof}
تعطي الخطية $\E(M_n) = \frac{n\mu}{n} = \mu$؛ ويعطي \omterm{def:g12:condprob:indep}{الاستقلال}
$\V(X_1 + \dots + X_n) = n\sigma^2$، والقسمة على $n$ تضرب \omterm{def:g12:randvar:exp}{التباين}
في $\frac{1}{n^2}$ (\cref{prop:g12:randvar:affine}).
\end{proof}

والمقدار $\frac{\sigma}{\sqrt n}$ هو \emph{قانون \omterm{def:g11:quad:discriminant}{الجذر} التربيعي} الأساسي: فلتنصيف
تذبذبات \omterm{def:g10:stats:mean}{متوسط}، ربّع حجم \omterm{def:g10:proba:sample}{العينة}.

\begin{omfigure}
\begin{tikzpicture}
\begin{axis}[omaxis,
  width=9.6cm, height=5.6cm,
  xlabel={$n$}, ylabel={$\sigma(M_n)/\sigma$},
  xmin=0, xmax=105, ymin=0, ymax=1.15,
  xtick={4,16,64}, ytick={1,0.5,0.25},
  yticklabels={$1$,$\tfrac12$,$\tfrac14$}]
  \addplot[omDef, thick, domain=1:100, samples=80] {1/sqrt(x)};
  \draw[black, dashed, thin] (axis cs:0,0.5) -- (axis cs:4,0.5)
    -- (axis cs:4,0);
  \draw[black, dashed, thin] (axis cs:0,0.25) -- (axis cs:16,0.25)
    -- (axis cs:16,0);
  \draw[black, dashed, thin] (axis cs:64,0.125) -- (axis cs:64,0);
  \fill (axis cs:1,1) circle (1.5pt);
  \fill (axis cs:4,0.5) circle (1.5pt);
  \fill (axis cs:16,0.25) circle (1.5pt);
  \fill (axis cs:64,0.125) circle (1.5pt);
\end{axis}
\end{tikzpicture}

{\small قانون \omterm{def:g11:quad:discriminant}{الجذر} التربيعي: كل تربيع لحجم \omterm{def:g10:proba:sample}{العينة} لا ينصّف
إلا \omterm{def:g11:stat:variance}{الانحراف المعياري} \omterm{def:g10:stats:mean}{للمتوسط}.}
\end{omfigure}

\section{متراجحات التركّز}

\begin{theorem}[متراجحة ماركوف]\label{thm:g12:sums:markov}
\index{متراجحة ماركوف}
إذا كان $X \geq 0$ و $a > 0$:
\[
\P(X \geq a) \leq \frac{\E(X)}{a}.
\]
\end{theorem}

\begin{proof}
في $\E(X) = \sum_i p_i x_i$ (وكل الحدود غير سالبة)، احتفظ فقط بالحدود
التي فيها $x_i \geq a$: فكل منها لا يقل عن $a\,p_i$، إذن
$\E(X) \geq a \sum_{x_i \geq a} p_i = a\,\P(X \geq a)$.
\end{proof}

\begin{theorem}[متراجحة بيينيمي--تشيبيشيف]\label{thm:g12:sums:chebyshev}
\index{متراجحة بيينيمي--تشيبيشيف}
من أجل أي \omterm{def:g12:randvar:rv}{متغير عشوائي} $X$ وأي $\varepsilon > 0$:
\[
\P\bigl(\abs{X - \E(X)} \geq \varepsilon\bigr)
\leq \frac{\V(X)}{\varepsilon^2}.
\]
\end{theorem}

\begin{proof}
طبّق متراجحة ماركوف على المتغير غير السالب
$Y = (X - \E(X))^2$ مع $a = \varepsilon^2$:
\[
\P\bigl(\abs{X - \E(X)} \geq \varepsilon\bigr)
= \P(Y \geq \varepsilon^2)
\leq \frac{\E(Y)}{\varepsilon^2} = \frac{\V(X)}{\varepsilon^2}. \qedhere
\]
\end{proof}

\begin{omfigure}
\begin{tikzpicture}
\begin{axis}[omaxis,
  width=9.6cm, height=5cm,
  ylabel={}, xmin=-3.4, xmax=3.4, ymin=0, ymax=1.25,
  xtick={-1,0,1},
  xticklabels={$\E(X)-\varepsilon$,$\E(X)$,$\E(X)+\varepsilon$},
  ytick=\empty]
  \addplot[name path=bell, omDef, thick, domain=-3.2:3.2]
    {exp(-x^2/2)};
  \path[name path=xaxis] (-3.2,0) -- (3.2,0);
  \addplot[omDef!20] fill between[of=bell and xaxis,
    soft clip={domain=-1:1}];
  \addplot[omThm!20] fill between[of=bell and xaxis,
    soft clip={domain=-3.2:-1}];
  \addplot[omThm!20] fill between[of=bell and xaxis,
    soft clip={domain=1:3.2}];
  \draw[black, dashed, thin] (axis cs:-1,0) -- (axis cs:-1,0.75);
  \draw[black, dashed, thin] (axis cs:1,0) -- (axis cs:1,0.75);
\end{axis}
\end{tikzpicture}

{\small التركّز: تحصر متراجحة بيينيمي--تشيبيشيف \omterm{def:g10:proba:distribution}{احتمال} أن
يقع $X$ في الذيول (بالأحمر)، على بعد $\varepsilon$ على الأقل من
أمله الرياضي، بالمقدار $\V(X)/\varepsilon^2$.}
\end{omfigure}

\begin{theorem}[قانون الأعداد الكبيرة]\label{thm:g12:sums:lln}
\index{قانون الأعداد الكبيرة}
لتكن $X_1, \dots, X_n$ مستقلة ولها \omterm{def:g12:randvar:rv}{التوزيع} نفسه
(\omterm{def:g12:randvar:exp}{الأمل الرياضي} $\mu$ \omterm{def:g12:randvar:exp}{والتباين} $\sigma^2$) وليكن $M_n$ \omterm{prop:g12:sums:mean}{متوسط عينتها}. عندئذٍ
من أجل كل $\varepsilon > 0$:
\[
\P\bigl(\abs{M_n - \mu} \geq \varepsilon\bigr)
\leq \frac{\sigma^2}{n\,\varepsilon^2}
\xrightarrow[n \to +\infty]{} 0 .
\]
فيتركّز \omterm{prop:g12:sums:mean}{متوسط العينة} حول \omterm{def:g12:randvar:exp}{الأمل الرياضي}.
\end{theorem}

\begin{proof}
متراجحة بيينيمي--تشيبيشيف مطبَّقة على $M_n$، الذي أمله الرياضي $\mu$ و
\omterm{def:g12:randvar:exp}{تباينه} $\frac{\sigma^2}{n}$ (\cref{prop:g12:sums:mean}).
\end{proof}

\begin{remark}
هذه المبرهنة جسر بين نظرية الاحتمالات والإحصاء: إذ إن \emph{\omterm{def:g10:stats:series}{تواتر}}
\omterm{def:g11:prob:model}{حادثة} على عدد كبير من التكرارات المستقلة يقترب من
\emph{احتمالها} (خذ $X_i$ مؤشر \omterm{def:g11:prob:model}{الحادثة}، فيكون
$\mu = p$). وهي تبرّر تقدير \omterm{def:g10:proba:distribution}{احتمال} بالمحاكاة
(\emph{طريقة مونت كارلو}) وتقدير نسبة في مجتمع باستطلاع ---
تقديرًا كميًا: انظر \cref{ch:g12:contdist}.
\end{remark}

\begin{method}[استعمال متراجحة بيينيمي--تشيبيشيف]\label{met:g12:sums:chebyshev}
لضمان $\P(\abs{M_n - \mu} \geq \varepsilon) \leq \alpha$، يكفي
أن يكون $n \geq \frac{\sigma^2}{\alpha\,\varepsilon^2}$. والحاصرة
خام (فالتذبذبات الحقيقية أصغر بكثير عادةً) لكنها عامة تمامًا:
إذ لا تقتضي شيئًا عن \omterm{def:g12:randvar:rv}{التوزيع} سوى \omterm{def:g12:randvar:exp}{تباين} منته.
\end{method}

\section{تمارين}

\begin{exercise}[$\star$]\label{exo:g12:sums:1}
يُرمى حجرا نرد غير مغشوشين؛ وليكن $S$ المجموع. باستعمال الخطية (لا
\omterm{def:g12:randvar:rv}{توزيع} $S$!)، احسب $\E(S)$؛ ثم احسب $\V(S)$ باستعمال
\omterm{def:g12:condprob:indep}{الاستقلال}، علمًا بأن \omterm{def:g12:randvar:exp}{تباين} حجر نرد واحد غير مغشوش هو $\frac{35}{12}$.
\end{exercise}

\begin{exercise}[$\star$]\label{exo:g12:sums:2}
ليكن $X \sim \mathcal B(100,\ 0.5)$. احصر
$\P(X \geq 75)$ بمتراجحة ماركوف، ثم
$\P(\abs{X - 50} \geq 25)$ بمتراجحة بيينيمي--تشيبيشيف. وقارن.
\end{exercise}

\begin{exercise}[$\star$]\label{exo:g12:sums:3}
تُرمى قطعة نقد غير مغشوشة $n$ مرة ويرمز $F_n$ إلى \omterm{def:g10:stats:series}{تواتر} الصور.
فكم يجب أن يكون $n$ كبيرًا حتى يكون، بحاصرة بيينيمي--تشيبيشيف،
$\P\bigl(\abs{F_n - 0.5} \geq 0.05\bigr) \leq 0.05$؟
\end{exercise}

\begin{exercise}[$\star\star$]\label{exo:g12:sums:4}
مستثمرة تقسم رأس مالها بالتساوي بين $n$ أصلًا مستقلًا،
لكل منها مردود مأمول $\mu = 5\%$ \omterm{def:g11:stat:variance}{وانحراف معياري}
$\sigma = 20\%$. احسب \omterm{def:g12:randvar:exp}{الأمل الرياضي} \omterm{def:g11:stat:variance}{والانحراف المعياري} لمردود
المحفظة $M_n$، وعدد الأصول اللازم لجعل
\omterm{def:g11:stat:variance}{الانحراف المعياري} دون $4\%$. وأي مبدأ مالي يوضّحه
هذا؟
\end{exercise}

\begin{exercise}[$\star\star$]\label{exo:g12:sums:5}
ليكن $X$ و $Y$ \omterm{def:g12:sums:indep}{مستقلين}، وكلاهما منتظم على $\{1, 2, 3\}$.
\begin{enumerate}
  \item أعطِ \omterm{def:g12:randvar:rv}{توزيع} $S = X + Y$ واحسب $\E(S)$ و $\V(S)$
        منه مباشرةً.
  \item استعد القيمتين بالخطية وبجمعية \omterm{def:g12:randvar:exp}{التباين}.
\end{enumerate}
\end{exercise}

\begin{exercise}[$\star\star$]\label{exo:g12:sums:6}
بيّن أن جمعية \omterm{def:g12:randvar:exp}{التباين} قد تسقط بلا \omterm{def:g12:condprob:indep}{استقلال}:
احسب $\V(X + Y)$ من أجل $Y = X$، وقارن بالمقدار $\V(X) + \V(Y)$. ومن أجل
أي متغيرات $X$ يكون $\V(2X) = 2\V(X)$؟
\end{exercise}

\begin{exercise}[$\star\star$]\label{exo:g12:sums:7}
يُشتبه في أن حجر نرد مغشوش. فيُرمى $1200$ مرة فيُظهر
ستة $260$ مرة ($f = 0.2167$ بدل $\frac16 \approx 0.1667$). وفي ظل
فرضية أن الحجر غير مغشوش، احصر
$\P\bigl(\abs{F_n - \frac16} \geq 0.05\bigr)$ بمتراجحة بيينيمي--تشيبيشيف، ثم
ناقش هل فرضية عدم الغش معقولة.
\end{exercise}

\begin{exercise}[$\star\star\star$]\label{exo:g12:sums:8}
\emph{(متراجحة أفضل من أجل قطعة النقد.)} ليكن
$X \sim \mathcal B(n,\ p)$ و $F_n = \frac Xn$.
\begin{enumerate}
  \item بيّن أن $p(1-p) \leq \frac14$ من أجل $p \in \intcc{0}{1}$.
  \item استنتج الحاصرة المستقلة عن \omterm{def:g12:randvar:rv}{التوزيع}
        $\P\bigl(\abs{F_n - p} \geq \varepsilon\bigr) \leq
        \frac{1}{4n\varepsilon^2}$.
  \item كم شخصًا يجب استطلاعه حتى يقع \omterm{def:g10:stats:series}{التواتر} الملاحَظ
        على بعد $3$ نقاط من النسبة الحقيقية \omterm{def:g10:proba:distribution}{باحتمال} لا يقل عن
        $95\%$، باستعمال هذه الحاصرة؟ (وتستعمل الاستطلاعات الحقيقية تقديرات أحدّ، لكن
        رتبة القدر صحيحة.)
\end{enumerate}
\end{exercise}

\section{مسألة: بيت اللعب يربح دائمًا}

\begin{problem}[{مسألة نهاية الأسبوع --- قانون الأعداد الكبيرة
يفسّر دور القمار والاستطلاعات والتأمين، ويهدم
مغالطة المقامر في الطريق}]
\label{pb:g12:sums:1}
لاعب روليت بعد $100$ دورة يكون، في أغلب الأحيان \omterm{def:g10:numbers:approx}{تقريبًا}،
\emph{رابحًا}. أما دار القمار التي تشغّل مليون دورة فرابحة
بيقين لا تشكّك فيه أي محكمة. اللعبة نفسها، والأفضلية الضئيلة نفسها ---
والفرق هو $\sqrt n$، وهو موضوع هذا
الفصل (\cref{prop:g12:sums:mean} و
\cref{thm:g12:sums:chebyshev} و \cref{thm:g12:sums:lln}). وتشغّل هذه
المسألة دفاتر دار القمار، وتحدّد حجم استطلاع انتخابي، وتسعّر
التنويع --- وتفكّك أغلى مغالطة في
تاريخ القمار.

\medskip
\noindent\textbf{الجزء الأول --- التمكّن.}
\begin{enumerate}
  \item حجرا نرد مستقلان: احسب $V(X + Y)$
        (\cref{thm:g12:sums:variance}).
  \item ارمِ $100$ حجر نرد مستقلًا وخذ \omterm{def:g10:stats:mean}{متوسط} النتائج:
        أعطِ \omterm{def:g12:randvar:exp}{الأمل الرياضي} \omterm{def:g11:stat:variance}{والانحراف المعياري} \omterm{prop:g12:sums:mean}{لمتوسط
        العينة} (\cref{prop:g12:sums:mean}؛ ومن أجل حجر واحد،
        $\sigma = \sqrt{35/12} \approx 1.71$).
  \item \omterm{def:g12:randvar:rv}{متغير عشوائي} غير سالب أمله الرياضي $2$.
        فماذا تقول متراجحة ماركوف
        (\cref{thm:g12:sums:markov}) عن
        $\P(X \geq 10)$؟
  \item احصر $\P\left(\abs{\bar X_{100} - 3.5} \geq
        0.5\right)$ من أجل \omterm{def:g10:stats:mean}{متوسط} حجرَي النرد في السؤال 2 بمتراجحة
        بيينيمي--تشيبيشيف.
  \item الآمال الرياضية تُجمع دائمًا، والتباينات لا تُجمع إلا عند
        \omterm{def:g12:condprob:indep}{الاستقلال}: أعطِ مثالًا لمتغيرين مرتبطين $X, Y$ (إرشاد:
        $Y = -X$) يكون من أجلهما
        $V(X + Y) \neq V(X) + V(Y)$.
\end{enumerate}

\medskip
\noindent\textbf{الجزء الثاني --- دفاتر بيت اللعب.}
الروليت الأوروبي، والرهان $1$ أورو على الأحمر: تربح $+1$
\omterm{def:g10:proba:distribution}{باحتمال} $\frac{18}{37}$، وتخسر $-1$ \omterm{def:g10:proba:distribution}{باحتمال}
$\frac{19}{37}$.
\begin{enumerate}[resume]
  \item احسب \omterm{def:g12:randvar:exp}{الأمل الرياضي} \omterm{def:g11:stat:variance}{والانحراف المعياري} لربح رهان
        واحد.
  \item مقامر يراهن $100$ رهانًا مستقلًا؛ وليكن $G$ الربح
        الكلي. احسب $\E(G)$ و $\sigma(G)$، ثم
        النسبة $\frac{\abs{\E(G)}}{\sigma(G)}$. فماذا يعني انزياح
        قدره ربع \omterm{def:g11:stat:variance}{انحراف معياري} بالنسبة إلى
        حظوظ المقامر في أن يكون رابحًا هذه الليلة؟
  \item ترى دار القمار $1\,000\,000$ رهان. احسب
        \omterm{def:g12:randvar:exp}{الأمل الرياضي} \omterm{def:g11:stat:variance}{والانحراف المعياري} لحصيلتها الكلية،
        والنسبة نفسها. وفسّر العدد $27$.
  \item صادِق بتشيبيشيف: احصر \omterm{def:g10:proba:distribution}{احتمال} أن \emph{تخسر}
        دار القمار مالًا على المليون رهان.
  \item أنظمة الرهان: المضاعفة بعد الخسائر، والانسحاب عند
        الربح\,\dots\ وباستعمال خطية \omterm{def:g12:randvar:exp}{الأمل الرياضي}
        (\cref{thm:g12:sums:linearity}) على \omterm{def:g12:seq:sequence}{متتالية} الرهانات (وقد
        تكون عشوائية)، فسّر لماذا يكون من أجل \emph{كل}
        استراتيجية في لعبة ذات أفضلية سالبة ربح مأمول سالب ---
        فماذا سينتهك نظام رابح؟
  \item صُغ ما يعد به قانون الأعداد الكبيرة
        (\cref{thm:g12:sums:lln}) بخصوص \omterm{def:g10:stats:series}{تواتر}
        الأحمر --- وما \emph{لا} يعد به بخصوص
        الدورة التالية بعد عشرة أحمر \omterm{def:g12:seq:sequence}{متتالية}. وسمِّ
        المغالطة.
  \item النقطة الأدق: القانون يعمل \emph{بالتخفيف}،
        لا بالتعويض. فإذا سبق الأحمر التوقع بمقدار $100$
        بعد أمسية ما، فلا يُردّ الفائض
        أبدًا --- احسب ما الذي يحدث بدلًا من ذلك \emph{للنسبة}
        $\frac{100}{n}$ عند $n = 10^6$، ثم
        أعد كتابة خطأ مغالطة المقامر في جملة واحدة.
\end{enumerate}

\medskip
\noindent\textbf{الجزء الثالث --- الاستطلاعات.}
\begin{enumerate}[resume]
  \item انطلاقًا من \cref{exo:g12:sums:8}: كم شخصًا يجب
        استطلاعه ليقع \omterm{def:g10:stats:series}{التواتر} الملاحَظ ضمن $3$
        نقاط من الحقيقة \omterm{def:g10:proba:distribution}{باحتمال} لا يقل عن
        $95\,\%$، بحاصرة تشيبيشيف المستقلة عن \omterm{def:g12:randvar:rv}{التوزيع}؟
  \item تستعمل معاهد الاستطلاع الحقيقية نحو $1\,100$ شخص من أجل
        هامش $\pm 3$ نقاط عند $95\,\%$: وصيغتهم هي
        $n \approx \frac{1.96^2 \times p(1-p)}
        {\varepsilon^2}$ حيث $p(1-p) \leq \frac14$. قدّرها،
        وفسّر الفجوة مع السؤال 13 (فما الذي
        لا تعرفه متراجحة تشيبيشيف عن شكل
        التذبذبات؟).
  \item لتنصيف هامش الخطأ، كيف يجب أن تنمو
        \omterm{def:g10:proba:sample}{العينة}؟ فمن $\pm 3$ نقاط عند $1\,100$، أي \omterm{def:g10:proba:sample}{عينة}
        تعطي $\pm 1.5$ نقطة؟
  \item الكلاسيكية المخالفة للحدس: لم يستعمل حجم \omterm{def:g10:proba:sample}{العينة} المطلوب
        حجم المجتمع أبدًا --- فإن $1\,100$ شخص
        يكفون لمدينة أو لقارة. أشر إلى الموضع
        في النموذج الذي يغيب فيه حجم المجتمع، و
        أعطِ تشبيه المطبخ الذي يستعمله المستطلعون.
\end{enumerate}

\medskip
\noindent\textbf{الجزء الرابع --- التنويع.}
\begin{enumerate}[resume]
  \item يكسب مؤمّن \cref{pb:g12:randvar:1} مبلغ
        $30n$ أورو أملًا رياضيًا على $n$ عقدًا \omterm{def:g11:stat:variance}{بانحراف
        معياري} للمطالبات يقارب $315\sqrt n$. فمن أجل
        أي $n$ يتجاوز الربح المأمول أخيرًا انحرافًا
        معياريًا واحدًا للمطالبات؟ وماذا يفعل
        $\sqrt n$ من أجل المؤمّن؟
  \item مردود سهم واحد السنوي يتذبذب بالمقدار
        $\sigma = 20\,\%$. قسّم المال بالتساوي على $25$
        سهمًا \emph{مستقلًا} كهذا: احسب
        $\sigma$ المحفظة. وتسمّي المالية التنويع
        الغداء المجاني الوحيد --- فما الغداء بالضبط؟
  \item والآن لتكن الأسهم $25$ مترابطة تمامًا (فتتحرك
        كلها معًا): فما $\sigma$ المحفظة؟
        قارن الطرفين وصُغ أي فرض تستأجره سرًا
        كل حجة تنويع --- وماذا حدث حين سقط على مستوى النظام كله سنة 2008.
  \item الخاتمة --- سيمفونية $\sqrt n$: المجاميع تتذبذب مثل
        $\sqrt n$ بينما تستقر متوسطاتها مثل
        $\frac{1}{\sqrt n}$؛ اعزف اللحن عبر الصناعات
        الأربع في هذه المسألة (دار القمار، والاستطلاع، والتأمين،
        والمحافظ)، وسمِّ المعتديَين القائمين (مغالطة
        المقامر \omterm{def:g12:condprob:indep}{والاستقلال} المزيف)، وأعطِ
        الإشارة إلى الأمام: \emph{شكل} التذبذبات
        --- أي الجرس --- هو نجم الفصل التالي
        المتصل، ومبرهنته الكاملة تتوّج الكتب
        الجامعية.
\end{enumerate}
\end{problem}
